\section{Related Work}
\label{sec:related_work}

\subsection{Evaluating Scientific Writing with LLMs}

% 科学写作评价正从通用文本相似度转向针对具体任务的质量判断。
% 早期研究主要通过与人类评审意见的重合度和作者感知的有用性
% 评价模型生成的论文反馈~\citep{liang-etal-2024-large-scale}。

Scientific-writing evaluation is moving beyond generic text similarity
toward task-specific quality judgments.
Early work evaluated generated paper feedback through overlap with human reviews
and author-perceived usefulness~\citep{liang-etal-2024-large-scale}.

% 近期方法使用任务特定的维度和标准评价不同的科学写作产物，
% 涵盖评审意见、相关工作和学术定位等任务
% ~\citep{sadallah-etal-2025-good,sahinuc-etal-2025-expert,
% xie-etal-2026-rwgbench}。在这类研究中，SurveyReview 将 rationale
% 质量与分数对齐分开评价~\citep{zhang-etal-2026-surveyreview}；
% Think-Probe-Respond 进一步发现，即使 rationale 接近专家推理，
% 最终判断仍可能存在偏差，说明这种区分具有必要性
% ~\citep{schopf-etal-2026-think-probe-respond}。
% 然而，rationale supervision 和推理协议如何分别影响
% rationale 质量与评分表现仍不清楚。

Recent work uses task-specific criteria to evaluate diverse
scientific-writing artifacts, including review comments, related work,
and scholarly positioning
~\citep{sadallah-etal-2025-good,sahinuc-etal-2025-expert,
xie-etal-2026-rwgbench}.
Within this line of work, SurveyReview evaluates rationale quality
separately from score alignment~\citep{zhang-etal-2026-surveyreview}.
Think-Probe-Respond further shows why this distinction matters:
even expert-like rationales can accompany miscalibrated final judgments
~\citep{schopf-etal-2026-think-probe-respond}.
However, how rationale supervision and inference protocols separately
affect rationale quality and scoring performance remains unclear.

% SciRM 更直接地体现了这一归因问题：它在动态标准和 rubric 下联合生成
% rationale 与分数，同时结合领域偏好优化、reasoning/self-verification
% 和 rationale-first 输出协议。因此，无法分别归因监督目标、
% 推理协议及其交互的影响~\citep{sahinuc-etal-2026-reward}。

SciRM more directly exposes this attribution problem: it jointly generates
rationales and scores under dynamic criteria and rubrics while combining
domain-preference optimization, reasoning/self-verification, and a
rationale-first output protocol. Because these components change together,
their effects cannot be attributed separately to the supervision target,
inference protocol, or their interaction~\citep{sahinuc-etal-2026-reward}

\subsection{Rationales in LLM Training and Inference}

% 在训练阶段，evaluator rationale 可以与分数共同作为监督目标，
% 但现有方法采用了不同的训练设计。Prometheus 联合学习评价反馈和分数；
% TRACT 则使用模型自生成的 reasoning trace 作为 rationale，
% 将预测分数替换为 gold score，并采用回归感知的分数学习目标
% ~\citep{kim-etal-2024-prometheus-iclr,chiang-etal-2025-tract}。
% 由于这些设计同时改变 rationale 来源和训练目标，
% rationale supervision 本身的影响难以单独识别。
% 这一归因问题也反映在不同的实验结果中：inferred thinking traces
% 可以提高 LLM rater 与人类判断的一致性
% ~\citep{zhang-etal-2025-judge-eyes}，但在临床预测中，
% 即使 rationale 准确且有证据支撑，rationale SFT 仍弱于
% label-only training~\citep{su-etal-2026-rationale-disease}。
% 因此，rationale supervision 在何种条件下有效仍不清楚。

At training time, evaluator rationales can be supervised alongside scores,
but existing methods implement this supervision differently.
Prometheus jointly learns evaluative feedback and scores, whereas TRACT uses
self-generated reasoning traces as rationales, replaces their predicted scores
with gold scores, and applies a regression-aware objective
~\citep{kim-etal-2024-prometheus-iclr,chiang-etal-2025-tract}.
Because these designs change the rationale source and training objective
together, the effect of rationale supervision is difficult to isolate.
This attribution problem is also reflected in mixed empirical results:
inferred thinking traces improve LLM--human agreement
~\citep{zhang-etal-2025-judge-eyes}, whereas rationale SFT underperforms
label-only training in clinical prediction even with accurate, grounded
rationales~\citep{su-etal-2026-rationale-disease}.
It therefore remains unclear when rationale supervision is beneficial.

% 在推理阶段，许多 LLM evaluator 会先生成 rationale，再给出分数。
% G-Eval 使用任务特定的评价步骤引导评分，并报告了更高的人类相关性
% ~\citep{liu-etal-2023-g}。但先生成分析再评分也可能压缩分数分布，
% 降低常用读出方式下的评价性能
% ~\citep{wang-etal-2025-improving-llm-judge}。
% Pointwise reranking 实验同样发现，即使采用针对性训练干预，
% CoT scoring 与 direct scoring 的性能差距仍然存在
% ~\citep{chen-etal-2026-beyond-polarization}。
% 因此，rationale-first inference 可能提高与人类判断的一致性，
% 也可能导致分数分布偏移；造成这些相反结果的原因仍不清楚。

At inference time, many LLM evaluators generate a rationale before assigning
a score. G-Eval uses task-specific evaluation steps and reports stronger
correlation with human judgments~\citep{liu-etal-2023-g}.
Conversely, generating an analysis before scoring can narrow the score
distribution and reduce performance under common readouts
~\citep{wang-etal-2025-improving-llm-judge}.
Pointwise-reranking experiments likewise find a persistent disadvantage for
CoT scoring relative to direct scoring despite targeted training interventions
~\citep{chen-etal-2026-beyond-polarization}.
Thus, rationale-first inference may improve human alignment but also distort
score distributions; the source of these conflicting effects remains unclear.

% 更广泛的 CoT 研究表明，显式推理的收益具有任务依赖性：
% 稳定增益主要集中在数学或符号推理任务
% ~\citep{sprague-etal-2025-cot}；对 LLM judge 的受控比较也发现，
% reasoning 有益于结构化验证，但在简单评价任务上的收益有限甚至为负
% ~\citep{zhang-etal-2026-reasoning-not-free}。
% 因此，有必要在科学写作评价中单独检验 rationale-based 方法。

Broader CoT studies show that explicit reasoning is task-dependent:
consistent gains concentrate mainly in mathematical or symbolic reasoning
~\citep{sprague-etal-2025-cot}, while controlled judge comparisons find
benefits for structured verification but limited or negative gains on simpler
evaluations~\citep{zhang-etal-2026-reasoning-not-free}.
These findings motivate a domain-specific examination of rationale-based
methods in scientific-writing evaluation.

% 现有受控研究已开始区分训练阶段的 rationale supervision 与推理阶段的
% reasoning。NLURC 比较 NLU 任务上的 label-only 与 rationale-augmented
% training，但未在两种输出接口下完整评估各训练方法
% ~\citep{shi-etal-2025-investigating}。TTCW 进一步交叉比较
% reasoning supervision 与 inference-time thinking，但研究的是长篇文学评论生成，
% 且性能下降主要来自解析失败~\citep{liu-etal-2026-when}。
% 因此，现有研究尚未在科学写作评分中完整交叉监督目标与推理协议，
% 也没有区分格式失败和合法分数中的评分错误。

Controlled studies have begun to separate training-time rationale supervision
from inference-time reasoning. NLURC compares label-only and
rationale-augmented training on NLU tasks but does not evaluate each method
under both output interfaces~\citep{shi-etal-2025-investigating}.
TTCW crosses reasoning supervision with inference-time thinking but studies
long-form literary review generation, where degradation is driven mainly by
parsing failures~\citep{liu-etal-2026-when}.
Existing work therefore neither fully crosses supervision targets with
inference protocols in scientific-writing scoring nor distinguishes format
failures from errors among valid scores.